%%%%%%%%%%%%%%%%
%%% num 13 %%%
%%%%%%%%%%%%%%%%

\Chapter{13}

\Verse{1}
{
	\hP{וַיְדַבֵּר יְהֹוָה אֶל מֹשֶׁה לֵּאמֹר׃}
}
{
	\eP{And YHWH spoke to Moses, saying,  }
}
{
}

\Verse{2}
{
	\hP{שְׁלַח לְךָ אֲנָשִׁים וְיָתֻרוּ אֶת אֶרֶץ כְּנַעַן אֲשֶׁר אֲנִי נֹתֵן לִבְנֵי יִשְׂרָאֵל אִישׁ אֶחָד אִישׁ אֶחָד לְמַטֵּה אֲבֹתָיו תִּשְׁלָחוּ כֹּל נָשִׂיא בָהֶם׃}
}
{
	\eP{
		“Send men and let them scout the land of Canaan that I’m
		giving to the children of Israel. You shall send one man
		for each tribe of his fathers, every one of them a
		chieftain.”
	}
}
{
%	\footnote{
%		send one man for each tribe. One scout is named for each tribe except
%		Levi. Why is the priestly tribe not represented? Perhaps no one from
%		Levi is sent to scout the territory because Levi is the one tribe that
%		will not be assigned any territory in the land. They derive their
%		support from the sacrifices, tithes, and donations of the other tribes.
%		Footnote 13:2. every one of them a chieftain. This story is commonly
%		referred to as the episode of the spies, but the men whom Moses sends
%		are tribal leaders, not ordinary scouts. The function of their mission
%		is not espionage but observation, to bring back a description of the
%		land to their people. They are leaders. That is why their negativity
%		when they return is so devastating. It is not just a disappointing bit
%		of intelligence. It is a failure to lead and encourage.
%	}
}

\Verse{3}
{
	\hP{וַיִּשְׁלַח אֹתָם מֹשֶׁה מִמִּדְבַּר פָּארָן עַל פִּי יְהֹוָה כֻּלָּם אֲנָשִׁים רָאשֵׁי בְנֵי יִשְׂרָאֵל הֵמָּה׃}
}
{
	\eP{
		And Moses sent them from the wilderness of Paran by YHWH’s
		word. They were all men who were heads of the children of
		Israel.
	}
}
{
}

\Verse{4}
{
	\hP{וְאֵלֶּה שְׁמוֹתָם לְמַטֵּה רְאוּבֵן שַׁמּוּעַ בֶּן זַכּוּר׃}
}
{
	\eP{
		And these are their names: For the tribe of Reuben,
		Shammua son of Zaccur.
	}
}
{
}

\Verse{5}
{
	\hP{לְמַטֵּה שִׁמְעוֹן שָׁפָט בֶּן חוֹרִי׃}
}
{
	\eP{For the tribe of Simeon, Shaphat son of Hori. }
}
{
}

\Verse{6}
{
	\hP{לְמַטֵּה יְהוּדָה כָּלֵב בֶּן יְפֻנֶּה׃}
}
{
	\eP{For the tribe of Judah, Caleb son of Jephunneh. }
}
{
}

\Verse{7}
{
	\hP{לְמַטֵּה יִשָּׂשכָר יִגְאָל בֶּן יוֹסֵף׃}
}
{
	\eP{For the tribe of Issachar, Igal son of Joseph. }
}
{
}

\Verse{8}
{
	\hP{לְמַטֵּה אֶפְרָיִם הוֹשֵׁעַ בִּן נוּן׃}
}
{
	\eP{For the tribe of Ephraim, Hoshea son of Nun. }
}
{
}

\Verse{9}
{
	\hP{לְמַטֵּה בִנְיָמִן פַּלְטִי בֶּן רָפוּא׃}
}
{
	\eP{For the tribe of Benjamin, Palti son of Raphu. }
}
{
}

\Verse{10}
{
	\hP{לְמַטֵּה זְבוּלֻן גַּדִּיאֵל בֶּן סוֹדִי׃}
}
{
	\eP{For the tribe of Zebulun, Gaddiel son of Sodi. }
}
{
}

\Verse{11}
{
	\hP{לְמַטֵּה יוֹסֵף לְמַטֵּה מְנַשֶּׁה גַּדִּי בֶּן סוּסִי׃}
}
{
	\eP{
		For the tribe of Joseph: of the tribe of Manasseh, Gaddi
		son of Susi.
	}
}
{
}

\Verse{12}
{
	\hP{לְמַטֵּה דָן עַמִּיאֵל בֶּן גְּמַלִּי׃}
}
{
	\eP{For the tribe of Dan, Ammiel son of Gemalli. }
}
{
}

\Verse{13}
{
	\hP{לְמַטֵּה אָשֵׁר סְתוּר בֶּן מִיכָאֵל׃}
}
{
	\eP{For the tribe of Asher, Sethur son of Michael. }
}
{
}

\Verse{14}
{
	\hP{לְמַטֵּה נַפְתָּלִי נַחְבִּי בֶּן ופְסִי׃}
}
{
	\eP{For the tribe of Naphtali, Nahbi son of Vophsi. }
}
{
}

\Verse{15}
{
	\hP{לְמַטֵּה גָד גְּאוּאֵל בֶּן מָכִי׃}
}
{
	\eP{For the tribe of Gad, Geuel son of Machi. }
}
{
}

\Verse{16}
{
	\hP{אֵלֶּה שְׁמוֹת הָאֲנָשִׁים אֲשֶׁר שָׁלַח מֹשֶׁה לָתוּר אֶת הָאָרֶץ וַיִּקְרָא מֹשֶׁה לְהוֹשֵׁעַ בִּן נוּן יְהוֹשֻׁעַ׃}
}
{
	\eP{
		These are the names of the men whom Moses sent to scout
		the land. And Moses called Hoshea son of Nun Joshua.
	}
}
{
}

\Verse{17}
{
	\hR{וַיִּשְׁלַח אֹתָם מֹשֶׁה לָתוּר אֶת אֶרֶץ כְּנָעַן}
	\hJ{וַיֹּאמֶר אֲלֵהֶם עֲלוּ זֶה בַּנֶּגֶב וַעֲלִיתֶם אֶת הָהָר׃}
}
{
	\eR{And Moses sent them to scout the land of Canaan.}
	\eJ{And he
		said to them, “Go up there in the Negeb, and you shall go
		up the mountain
	}
}
{
}

\Verse{18}
{
	\hJ{וּרְאִיתֶם אֶת הָאָרֶץ מַה הִוא וְאֶת הָעָם הַיֹּשֵׁב עָלֶיהָ הֶחָזָק הוּא הֲרָפֶה הַמְעַט הוּא אִם רָב׃}
}
{
	\eJ{
		and see the land, how it is; and the people who live on
		it, are they strong or weak, are they few or many;
	}
}
{
}

\Verse{19}
{
	\hJ{וּמָה הָאָרֶץ אֲשֶׁר הוּא יֹשֵׁב בָּהּ הֲטוֹבָה הִוא אִם רָעָה וּמָה הֶעָרִים אֲשֶׁר הוּא יוֹשֵׁב בָּהֵנָּה הַבְּמַחֲנִים אִם בְּמִבְצָרִים׃}
}
{
	\eJ{
		and how is the land in which they live, is it good or bad;
		and how are the cities in which they live, are they in
		camps or in fortified places;
	}
}
{
}

\Verse{20}
{
	\hJ{וּמָה הָאָרֶץ הַשְּׁמֵנָה הִוא אִם רָזָה הֲיֵשׁ בָּהּ עֵץ אִם אַיִן וְהִתְחַזַּקְתֶּם וּלְקַחְתֶּם מִפְּרִי הָאָרֶץ וְהַיָּמִים יְמֵי בִּכּוּרֵי עֲנָבִים׃}
}
{
	\eJ{
		and how is the land, is it fat or meager; does it have
		trees or not; and exert strength and take some fruit of
		the land.” And it was the days of the first grapes.
	}
}
{
}

\Verse{21}
{
	\hP{וַיַּעֲלוּ וַיָּתֻרוּ אֶת הָאָרֶץ מִמִּדְבַּר צִן עַד רְחֹב לְבֹא חֲמָת׃}
}
{
	\eP{
		And they went up and scouted the land from the wilderness
		of Zin to Rehob at the entrance of Hamath.
	}
}
{
}

\Verse{22}
{
	\hP{וַיַּעֲלוּ בַנֶּגֶב וַיָּבֹא עַד חֶבְרוֹן וְשָׁם אֲחִימַן שֵׁשַׁי וְתַלְמַי יְלִידֵי הָעֲנָק וְחֶבְרוֹן שֶׁבַע שָׁנִים נִבְנְתָה לִפְנֵי צֹעַן מִצְרָיִם׃}
}
{
	\eP{
		And they went up in the Negeb and came to Hebron. And
		Ahiman, Sheshai, and Talmai, the offspring of the giants,
		were there. (And Hebron was built seven years before Zoan
		of Egypt.)
	}
}
{
%	\footnote{
%		came to Hebron. They return to a place where Abraham lived (Gen 13:18),
%		where he met the three visitors and first questioned God (Genesis 18;
%		Mamre is at Hebron), and where the patriarchs and matriarchs are buried
%		(Gen 23:19). It is also the place from which Jacob sent Joseph to find
%		his brothers, which, in a sense, initiated the sequence of events that
%		led to Israel’s being in Egypt. It is therefore, meaningful in a variety
%		of ways that the scouts are especially noted to come to Hebron. It
%		particularly symbolizes that this is a return home. It might have been
%		the perfect place to begin Israel’s return to the land, but the people’s
%		resistance will change all that.
%	}
}

\Verse{23}
{
	\hJ{וַיָּבֹאוּ עַד נַחַל אֶשְׁכֹּל וַיִּכְרְתוּ מִשָּׁם זְמוֹרָה וְאֶשְׁכּוֹל עֲנָבִים אֶחָד וַיִּשָּׂאֻהוּ בַמּוֹט בִּשְׁנָיִם וּמִן הָרִמֹּנִים וּמִן הַתְּאֵנִים׃}
}
{
	\eJ{
		And they came to the Wadi Eshcol, and they cut a branch
		and one cluster of grapes from there—and carried it on a
		pole by two people—and some pomegranates and some figs.
	}
}
{
}

\Verse{24}
{
	\hJ{לַמָּקוֹם הַהוּא קָרָא נַחַל אֶשְׁכּוֹל עַל אֹדוֹת הָאֶשְׁכּוֹל אֲשֶׁר כָּרְתוּ מִשָּׁם בְּנֵי יִשְׂרָאֵל׃}
}
{
	\eJ{
		That place was called Wadi Eshcol on account of the
		cluster that the children of Israel cut from there.
	}
}
{
}

\Verse{25}
{
	\hP{וַיָּשֻׁבוּ מִתּוּר הָאָרֶץ מִקֵּץ אַרְבָּעִים יוֹם׃}
}
{
	\eP{
		And they came back from scouting the land at the end of
		forty days.
	}
}
{
}

\Verse{26}
{
	\hP{וַיֵּלְכוּ וַיָּבֹאוּ אֶל מֹשֶׁה וְאֶל אַהֲרֹן וְאֶל כּל עֲדַת בְּנֵי יִשְׂרָאֵל אֶל מִדְבַּר פָּארָן קָדֵשָׁה וַיָּשִׁיבוּ אֹתָם דָּבָר וְאֶת כּל הָעֵדָה וַיַּרְאוּם אֶת פְּרִי הָאָרֶץ׃}
}
{
	\eP{
		And they went and came to Moses and to Aaron and to all
		the congregation of the children of Israel, to the
		wilderness of Paran, at Kadesh; and they brought back word
		to them and all the congregation and showed them the
		land’s fruit.
	}
}
{
}

\Verse{27}
{
	\hJ{וַיְסַפְּרוּ לוֹ וַיֹּאמְרוּ בָּאנוּ אֶל הָאָרֶץ אֲשֶׁר שְׁלַחְתָּנוּ וְגַם זָבַת חָלָב וּדְבַשׁ הִוא וְזֶה פִּרְיָהּ׃}
}
{
	\eJ{
		And they told him and said, “We came to the land where you
		sent us, and also it’s flowing with milk and honey, and
		this is its fruit.
	}
}
{
}

\Verse{28}
{
	\hJ{אֶפֶס כִּי עַז הָעָם הַיֹּשֵׁב בָּאָרֶץ וְהֶעָרִים בְּצֻרוֹת גְּדֹלֹת מְאֹד וְגַם יְלִדֵי הָעֲנָק רָאִינוּ שָׁם׃}
}
{
	\eJ{
		Nonetheless: the people who live in the land are strong.
		And the cities are fortified, very big. And also we saw
		the offspring of the giants there.
	}
}
{
}

\Verse{29}
{
	\hJ{עֲמָלֵק יוֹשֵׁב בְּאֶרֶץ הַנֶּגֶב וְהַחִתִּי וְהַיְבוּסִי וְהָאֱמֹרִי יוֹשֵׁב בָּהָר וְהַכְּנַעֲנִי יוֹשֵׁב עַל הַיָּם וְעַל יַד הַיַּרְדֵּן׃}
}
{
	\eJ{
		Amalek lives in the land of the Negeb, and the Hittite and
		the Jebusite and the Amorite live in the mountains, and
		the Canaanite lives by the sea and along the Jordan.”
	}
}
{
}

\Verse{30}
{
	\hJ{וַיַּהַס כָּלֵב אֶת הָעָם אֶל מֹשֶׁה וַיֹּאמֶר עָלֹה נַעֲלֶה וְיָרַשְׁנוּ אֹתָהּ כִּי יָכוֹל נוּכַל לָהּ׃}
}
{
	\eJ{
		And Caleb quieted the people toward Moses and said, “Let’s
		go up, and we’ll take possession of it, because we’ll be
		able to handle it.”
	}
}
{
}

\Verse{31}
{
	\hJ{וְהָאֲנָשִׁים אֲשֶׁר עָלוּ עִמּוֹ אָמְרוּ לֹא נוּכַל לַעֲלוֹת אֶל הָעָם כִּי חָזָק הוּא מִמֶּנּוּ׃}
}
{
	\eJ{
		And the men who went up with him said, “We won’t be able
		to go up against the people, because they’re stronger than
		we are.”
	}
}
{
}

\Verse{32}
{
	\hP{וַיֹּצִיאוּ דִּבַּת הָאָרֶץ אֲשֶׁר תָּרוּ אֹתָהּ אֶל בְּנֵי יִשְׂרָאֵל לֵאמֹר הָאָרֶץ אֲשֶׁר עָבַרְנוּ בָהּ לָתוּר אֹתָהּ אֶרֶץ אֹכֶלֶת יוֹשְׁבֶיהָ הִוא וְכל הָעָם אֲשֶׁר רָאִינוּ בְתוֹכָהּ אַנְשֵׁי מִדּוֹת׃}
}
{
	\eP{
		And they brought out a report of the land that they had
		scouted to the children of Israel, saying, “The land
		through which we passed to scout it: it’s a land that eats
		those who live in it, and all the people whom we saw in it
		were people of size!
	}
}
{
}

\Verse{33}
{
	\hJ{וְשָׁם רָאִינוּ אֶת הַנְּפִילִים בְּנֵי עֲנָק מִן הַנְּפִלִים וַנְּהִי בְעֵינֵינוּ כַּחֲגָבִים וְכֵן הָיִינוּ בְּעֵינֵיהֶם׃}
}
{
	\eJ{
		And we saw the Nephilim there, sons of giants from the
		Nephilim, and we were like grasshoppers in our eyes, and
		so were we in their eyes.”
	}
}
{
%	\footnote{
%		sons of giants from the Nephilim. The huge creatures known as the
%		Nephilim come up in a chain of references that are spread far apart
%		through several books of the Tanak. Their origin is explained near the
%		beginning of the Torah, in the story of the “sons of God” who have
%		relations with human women who then give birth to giants (Gen 6:1–4).
%		Now in the spies episode, the Israelite spies see the descendants of
%		these giants (Numbers 13). Later, Joshua eliminates the giants from all
%		of the land except from the Philistine cities, including Gath (Josh
%		11:21–22). And later still, the Philistine giant Goliath comes from Gath
%		(1 Sam 17:4). These widespread references are among the many stitches
%		that bind the Tanak together as a continuing story. Those who see the
%		Bible as a loose collection of stories fail to take proper account of
%		the substantial number of such widely distributed binding elements.
%	}
}
